\section{PR-II Exact Classifications and Surviving Branches}
\label{sec:pr2-branch-classification}

PR-II extends the PR-I foundation to an electroweak and flavor architecture \cite{PRII}.  The frozen-postulate test separates exact consequences of a fixed electroweak representation from choices not selected by PR-I.

\subsection{Hypercharge Solution within Fixed Matter Content}

Write the one-generation left-handed Weyl hypercharges \cite{Weinberg_QFT_1996} as $q,u,d,\ell,e$ for the quark doublet, up conjugate, down conjugate, lepton doublet, and charged-lepton conjugate.  With $q=t$, the linear local-anomaly equations reduce to
\begin{equation}
 \ell=-3t,\qquad e=6t,\qquad d=-2t-u.
\end{equation}
Substitution into the cubic anomaly gives the exact factorization
\begin{equation}
 -18t(u-2t)(u+4t)=0.
 \label{eq:pr35-hypercharge-factorization}
\end{equation}
For $t\ne0$, the last two factors yield the standard branch and its up/down singlet relabeling.  Overall hypercharge normalization is conventional.  The $t=0$ factor gives an anomaly-free vectorlike singlet family $q=\ell=e=0$, $d=-u$; it is excluded only when the nonzero residual-charge ledger is imposed.  Hence anomaly cancellation alone is not a unique derivation, whereas anomaly cancellation plus a fixed nonzero charge ledger does isolate the familiar branch modulo normalization and labeling.

\subsection{Unique Photon Direction Conditional on the Order Representation}

Once the neutral mass matrix is derived in outer-product form,
\begin{equation}
 M_0^2\propto
 \begin{pmatrix}g^2&-gg'\\-gg'&g'^2\end{pmatrix}
 =\begin{pmatrix}g\\-g'\end{pmatrix}
  \begin{pmatrix}g&-g'\end{pmatrix},
\end{equation}
nonzero $g,g'$ imply rank one and a one-dimensional nullspace.  The photon direction is unique up to normalization.  This result does not derive the outer-product form from P1--P7; determinant zero by itself admits infinitely many symmetric rank-one matrices.

\subsection{Generation Replication is an Exact Counterexample}

Let one complete anomaly-free generation contain four weak doublets after color multiplicity.  Replicating the complete ledger $N$ times multiplies every local anomaly coefficient by $N$, so all remain zero.  The number of weak doublets is $4N$, which is even for every $N\in\mathbb Z_{>0}$; the Witten \cite{Witten_1982} parity condition therefore also holds.  Consequently,
\begin{equation}
 N=1,2,3,\ldots
\end{equation}
is an infinite discrete family satisfying these gates.  Neither anomaly cancellation nor Witten parity selects $N_{\rm gen}=3$.  Identifying the dimension of the broken quotient with a generation count is an extra selector unless an index theorem from the frozen PR-I postulates proves that identification.  The PR-II flavor, hierarchy, and neutrino outputs that depend on three sheets cannot therefore enter the unconditional closed theorem.

The same conclusion applies upstream: the electroweak lift $SU(2)_L\times U(1)_Y$, the weak boundary-deficit rule, hierarchy action, flavor overlap rules, neutrino branch, and strong anchor do not appear as consequences of P1 through P7 in the public derivation graph.  They remain conditional structures, not new postulates.
